% 自测题3计算3.3.1—3.3.7与问答3.4.1；
% 原书印刷56—57页／PDF物理64—65页，特别注明同时性与能量定义。
\subsection{自测题3：洛伦兹变换、时间间隔和质能}

\sourcetest{3.3.1}{56}{64}
\problemtext 甲、乙静止于惯性系$K,K'$，
甲测得两个同地事件的时间间隔为4s，
乙测得对应间隔为5s。
求$K'$相对$K$的速率，及乙测得这两个事件的空间距离。
\workings
甲系两事件同地，故4s为固有时间。
$\gamma=5/4$，于是
\[
\boxed{v=c\sqrt{1-(4/5)^2}=0.60c}.
\]
若$K'$沿$+x$相对$K$移动，
$\Delta x=0$且$\Delta t=4$s，
由洛伦兹变换
\[
\boxed{\Delta x'=-\gamma v\Delta t
=-(5/4)(0.6c)(4\,\mathrm s)
=-3c\,\mathrm s\simeq-9.0\times10^8\,\mathrm m}.
\]
空间\emph{间隔大小}为$9.0\times10^8$m。
\crosscheck 原题只给相对速率大小，
因此有向位移正负依选定坐标轴；不能仅依据两时间读出唯一方向。

\sourcetest{3.3.2}{57}{65}
\problemtext 飞船固有长度$L_0=90$m，
对地速度$0.8c$，掠过地面一固定观测站。
求(1)观测站地面时钟记录整艘飞船通过的时间，
(2)飞船本地系测得观测站从船首到船尾经过的时间。
\workings
地面同时测长$L=L_0\sqrt{1-0.8^2}=54$m，
所以地面观测站上同地两事件所隔的固有时间
\[
\boxed{\Delta t_{\rm ground}=
\frac L{0.8c}=\frac{54}{0.8c}
\simeq2.25\times10^{-7}\,\mathrm s}.
\]
飞船系中本船静止长度为90m，
地面站以$0.8c$穿过两端，故
\[
\boxed{\Delta t_{\rm ship}=
\frac{L_0}{0.8c}
\simeq3.75\times10^{-7}\,\mathrm s}.
\]
\crosscheck 船首遇站和船尾遇站在\emph{地面同地}、
在\emph{飞船异地}，因此两系时间间隔不相同；
$\Delta t_{\rm ship}=\gamma\Delta t_{\rm ground}$。

\sourcetest{3.3.3}{57}{65}
\problemtext 在$K$系中两个事件沿$x$的间距
$\Delta x=500$m、时间差$\Delta t=2\times10^{-7}$s。
若$K'$系观察者测得这两个事件同时，
求$K'$沿$x$相对$K$的有向速度。
\workings
$\Delta t'=\gamma(\Delta t-v\Delta x/c^2)=0$，
从而
\[
\boxed{v=\frac{c^2\Delta t}{\Delta x}
=\frac{(3.0\times10^8)^2(2\times10^{-7})}{500}
=3.6\times10^7\,\mathrm{m/s}=0.12c}.
\]
\crosscheck $|\Delta x|>c|\Delta t|$，事件为空间型间隔，
存在将其视为同时的惯性系；若是类时分离则不可如此变换。

\sourcetest{3.3.4}{57}{65}
\problemtext 地面发射台左右的西站W、东站E距台均为$L_0$。
发射台同一时刻向两站发出同一光信号。
一架飞机以$v$沿从W到E的方向匀速飞过，
求飞机上记录两接收站\emph{收到同一信号}
这两个事件的时间间隔。
\workings
发射台与两站均静止于地面系$K$。
两接收事件在地面系同时：
$t_E=t_W=L_0/c$，位置
$x_E=+L_0$、$x_W=-L_0$，
所以$\Delta t=0$、$\Delta x=2L_0$。
飞机系沿$+x$相对地面以$v$运动，
\[
\Delta t'=t_E'-t_W'
=\gamma\left(\Delta t-\frac{v\Delta x}{c^2}\right)
=-\frac{2\gamma vL_0}{c^2}.
\]
故
\[
\boxed{|\Delta t'|=
\frac{2vL_0}{c^2\sqrt{1-v^2/c^2}}}.
\]
负号表明飞机系认为\textbf{东站先接收}。
\crosscheck 切不可把两接收站误当成\emph{随飞机运动}，
也不可用$L_0/(c-v)$描述地面静止站的接收时刻。

\sourcetest{3.3.5}{57}{65}
\problemtext 一宇宙射线介子动能
$K=7M_0c^2$，
其中$M_0$是其静止质量，
求实验室系测得寿命与介子固有平均寿命之比。
\workings
$K=(\gamma-1)M_0c^2=7M_0c^2$，
故$\gamma=8$。
时间延缓使
\[
\boxed{\tau_{\rm lab}/\tau_0=8}.
\]
\crosscheck 不可直接把动能是7倍静能误写为时间延缓因子7；
延缓因子是总能量与静能之比8。

\sourcetest{3.3.6}{57}{65}
\problemtext 将电子速率从$1.2\times10^8$m/s
提高至$2.4\times10^8$m/s，
求所需外力的相对论净功，
取电子静质量$9.11\times10^{-31}$kg。
\workings
采用$c\simeq3.0\times10^8$m/s，
$\beta_1=0.4,\beta_2=0.8$，
$\gamma_1=1/\sqrt{0.84}$、
$\gamma_2=1/\sqrt{0.36}=5/3$。
相对论功—能关系
\[
\begin{aligned}
W&=(\gamma_2-\gamma_1)m_0c^2\\
&=\left(\frac53-\frac1{\sqrt{0.84}}\right)
(9.11\times10^{-31})(3.0\times10^8)^2,
\end{aligned}
\]
故
\[
\boxed{W\simeq4.72\times10^{-14}\,\mathrm J
\simeq0.294\,\mathrm{MeV}}.
\]
\crosscheck 只改变动能而不改变静能，
故总能量差与动能差相同。

\sourcetest{3.3.7}{57}{65}
\problemtext 观察者甲以$0.8c$相对静止的乙运动，
甲携带静止质量1kg的物体。
两人各自测得该物体的\emph{总能量}是多少？
\workings
物体在甲系静止，所以
\[
\boxed{E_{\rm A}=mc^2
\simeq8.99\times10^{16}\,\mathrm J}.
\]
乙系中物体以$0.8c$运动，
$\gamma=1/\sqrt{1-0.8^2}=5/3$，
\[
\boxed{E_{\rm B}=\gamma mc^2
=\tfrac53mc^2
\simeq1.50\times10^{17}\,\mathrm J}.
\]
\crosscheck 总能量不是洛伦兹标量，
\emph{静质量}才是；甲不能看到自己携带的物体有额外平动动能。

\subsection{自测题3：经典与相对论不变量问答}

\sourcetest{3.4.1}{57}{65}
\problemtext 对位移、质量、时间、速度、动能这五个物理量，
比较其在经典伽利略变换与狭义相对论洛伦兹变换下的表达，
指出哪些可保持不变。
\workings
一维伽利略变换为
$x'=x-vt,\ t'=t$，速度$u'=u-v$；
静质量$m'=m$、经典动能
$K'=\tfrac12m(u-v)^2$，
故\textbf{时间差和质量}一般保持不变；
\emph{同一时刻两点的空间距离}也保持不变，
但不同时间发生的两事件位移差
$\Delta x'=\Delta x-v\Delta t$并非一般不变量。
洛伦兹变换为
\[
x'=\gamma(x-vt),\qquad
t'=\gamma(t-vx/c^2).
\]
长度/位移、坐标时间、速度、动能一般随惯性系改变，
现代定义的\textbf{静质量$m_0$保持不变}，
同时有时空间隔
\[
\boxed{c^2\Delta t'^2-\Delta x'^2
=c^2\Delta t^2-\Delta x^2}
\]
和单粒子质能关系
$\boxed{E^2-(pc)^2=m_0^2c^4}$保持不变。
\crosscheck 若“质量”沿用原书旧式\emph{运动质量}
$m_{\rm rel}=\gamma m_0$，则它随速度变，不是洛伦兹标量；
因此应事先说明所说的是现代\emph{静质量}
还是教材历史术语的\emph{运动质量}。
